
\chapter{Background and Related Work}
\label{ch:related-work}

This chapter surveys the literatures that converge on the problem
addressed in this dissertation: battery dispatch optimisation, telemetry
faults in cyber-physical systems, uncertainty quantification via conformal
methods, and runtime safety assurance.  For each strand, we identify the
specific gap that the ORIUS framework fills.

\section{Battery Dispatch and Renewable-Aware Control}

Grid-scale battery dispatch is the problem of choosing charge and
discharge actions to maximise economic value (arbitrage, ancillary
services) while respecting SOC, power, and ramp constraints.  The
literature falls into three broad families.

\subsection{Deterministic Dispatch}

Linear programming (LP) formulations treat load, solar, and wind
forecasts as point estimates and solve a single deterministic optimisation
at each decision epoch~\cite{xu2020optimal}.  The approach is fast and
widely deployed, but it ignores forecast uncertainty entirely: the action
is optimal only if the forecast is exact.  Under telemetry degradation,
the deterministic LP has no mechanism to detect that its inputs are stale
or incorrect, and it will dispatch against an observation that may not
reflect the plant state.

\subsection{Robust and Tube-Based Methods}

Robust MPC constructs a \emph{tube} of feasible trajectories around the
nominal forecast, ensuring that the dispatched action is safe for all
realisations within a bounded uncertainty set~\cite{bemporad2007robust}.
Tube-based approaches are conservative by construction: they sacrifice
economic performance to guarantee constraint satisfaction.  However, the
uncertainty set is typically derived from forecast error statistics and
does not account for observation-quality degradation.  If the tube is
calibrated on clean telemetry, it will be \emph{too narrow} when telemetry
is degraded---precisely the regime where safety margins matter most.

\subsection{Stochastic and Risk-Aware Methods}

Stochastic programming and chance-constrained formulations model
uncertainty through scenario trees or distributional
assumptions~\cite{shapiro2021lectures}.  Risk-aware variants incorporate
CVaR or other coherent risk measures to penalise tail outcomes.  These
methods acknowledge uncertainty but assume it is \emph{aleatoric}
(inherent randomness in load and renewables) rather than \emph{epistemic}
(degraded or missing observations).  None of the existing risk-aware
dispatch frameworks condition their uncertainty sets on the quality of the
incoming telemetry.

\section{The Measurement Blind Spot}

Across all three families, a common blind spot emerges: the optimiser
trusts its inputs.  Whether the formulation is deterministic, robust, or
stochastic, it assumes that the observation~$\hat{x}_t$ arriving at the
decision epoch is a faithful---if noisy---representation of the true
state~$x_t$.  This assumption is codified in the measurement
equation~$\hat{x}_t = x_t + v_t$ with well-characterised noise~$v_t$.

In practice, telemetry faults violate this equation structurally.  A
dropped packet does not add noise; it removes information.  A stale
sensor does not perturb the reading; it freezes it.  A replay attack
substitutes a plausible but outdated value.  These faults create
\emph{epistemic} degradation that standard noise models cannot capture.

The measurement blind spot is the root cause of the Observation-Action
Safety Gap defined in Chapter~\ref{ch:introduction}: the optimiser
dispatches against a state that may not exist, and no downstream
constraint check can detect the error because it, too, operates on the
same degraded observation.

\section{Telemetry Faults in Cyber-Physical Systems}

The CPS literature documents telemetry faults extensively in the context
of power systems~\cite{liang2017review}, industrial
SCADA~\cite{cardenas2011attacks}, and autonomous
vehicles~\cite{sun2020adversarial}.  Fault taxonomies typically
distinguish:

\begin{description}
  \item[Packet dropout:] The reading is absent; the system either holds
    the last known value or interpolates.
  \item[Delay and jitter:] The reading arrives late; the decision is made
    on stale information.
  \item[Spike:] A transient hardware glitch produces an out-of-range
    value.
  \item[Stale sensor:] The sensor output is constant across multiple
    intervals, indicating a frozen transducer or a repeated value.
  \item[Covariate drift:] The joint distribution of load, solar, and wind
    shifts between the calibration period and deployment, invalidating
    learned residual distributions.
\end{description}

Existing fault-tolerant control methods address each fault type in
isolation---median filters for spikes, watchdog timers for stale
sensors, Kalman-based imputation for dropouts.  None synthesise the
fault signals into a \emph{unified reliability score} that feeds back
into the safety layer.  The Observation Quality Engine (OQE) introduced
in this dissertation fills this gap.

\section{Uncertainty Quantification and Conformal Methods}

\subsection{Conformal Quantile Regression (CQR)}

Conformal prediction provides distribution-free prediction intervals with
finite-sample coverage guarantees~\cite{vovk2005algorithmic}.  Conformal
Quantile Regression (CQR) combines quantile regression with conformal
calibration to produce adaptive intervals whose width varies with the
input features~\cite{romano2019conformalized}.

CQR is attractive for battery dispatch because it makes no distributional
assumptions about the forecast residuals and inherits the exchangeability
guarantee of conformal inference.  The coverage guarantee states:
\[
  \Pr\!\bigl[\,Y_{n+1} \in \hat{C}(X_{n+1})\,\bigr]
  \;\ge\; 1 - \alpha,
\]
where $\hat{C}$ is the conformal prediction set and $\alpha$ is the
target miscoverage level.

\subsection{Adaptive Conformal Inference}

Adaptive conformal inference (ACI) adjusts the conformal threshold online
to maintain coverage under distribution shift~\cite{gibbs2021adaptive}.
ACI is essential for battery dispatch, where load and renewable patterns
exhibit diurnal, weekly, and seasonal non-stationarity.

\subsection{Key Limitation}

The standard CQR/ACI guarantee assumes that calibration residuals are
exchangeable with test residuals.  Under telemetry degradation, this
assumption may be violated: a dropout burst changes the residual
distribution.  The ORIUS framework addresses this limitation by
\emph{inflating} the conformal interval width in proportion to the
observation quality penalty~$1/w_t$, restoring coverage even when the
exchangeability assumption is weakened.

\section{Runtime Safety, Shields, and Assurance}

Runtime verification and shielding have a rich history in formal methods
and reactive synthesis~\cite{bloem2015shield}.  A \emph{shield} is a
runtime monitor that intercepts unsafe actions before they reach the
plant, replacing them with the nearest safe alternative.

\subsection{How DC3S Extends the State of the Art}

Existing shields operate on the \emph{observed} state: they check whether
the proposed action satisfies constraints given the current observation.
DC3S introduces three innovations:

\begin{enumerate}
  \item \textbf{Observation-quality conditioning.}  The shield's
    constraint tightening is a function of the OQE score~$w_t$, not a
    fixed margin.  When telemetry is clean ($w_t \approx 1$), the shield
    is nearly transparent; when telemetry degrades, the shield tightens
    progressively.
  \item \textbf{Conformal uncertainty integration.}  The tightened
    constraints are derived from CQR intervals inflated by $1/w_t$,
    inheriting the conformal coverage guarantee.
  \item \textbf{Certified actions.}  Every dispatched action carries a
    machine-readable certificate recording the observation quality,
    uncertainty bounds, and safety margin, enabling downstream audit
    and temporal-validity analysis
    (Chapter~\ref{ch:cert-half-life-blackout}).
\end{enumerate}

\section{Control Barrier Functions}

Control barrier functions (CBFs) provide a Lyapunov-like certificate for
forward invariance of a safe set~\cite{ames2019control}.  A CBF
$h : \mathcal{X} \to \mathbb{R}$ defines the safe set as
$\mathcal{C} = \{x : h(x) \ge 0\}$ and enforces the condition
$\dot{h}(x, u) + \gamma\, h(x) \ge 0$ for some class-$\mathcal{K}$
function~$\gamma$.

CBFs assume \emph{full state feedback}: the barrier condition is evaluated
at the true state~$x_t$.  When only a degraded observation~$\hat{x}_t$ is
available, the CBF condition may be satisfied on the observation while
violated on the truth---precisely the OASG.  DC3S can be viewed as a
\emph{data-centric CBF} that inflates the barrier margin based on
observation quality, closing the gap that classical CBFs leave open.

\section{Information-Theoretic Roots}

The observation quality score~$w_t$ can be interpreted through an
information-theoretic lens.  When telemetry is perfect, the mutual
information $I(X_t; \hat{X}_t) = H(X_t)$ (full information).  Each fault
type reduces $I(X_t; \hat{X}_t)$: dropout removes entire samples, delay
introduces conditional-entropy growth, and stale sensors reduce the
effective channel capacity to zero.

The OQE multiplicative formula
\[
  w_t = \max\!\bigl(w_{\min},\;
  p_{\mathrm{miss}} \cdot p_{\mathrm{delay}} \cdot p_{\mathrm{ooo}} \cdot p_{\mathrm{spike}}\bigr)
\]
is a heuristic proxy for the normalised mutual information
$I(X_t; \hat{X}_t) / H(X_t)$.  While not a tight bound, the
multiplicative structure ensures that any single fault channel can
independently drive reliability toward the floor~$w_{\min}$, matching
the information-theoretic intuition that a single broken channel
dominates the overall quality.

\section{Evaluation Protocols and the Benchmark Gap}

Standard battery dispatch benchmarks evaluate controllers on economic
metrics (revenue, regret) and observed constraint satisfaction.  None
maintain a parallel truth trajectory, so they \emph{cannot detect the
OASG} by construction.  The CPSBench harness introduced in this
dissertation (Chapter~\ref{ch:orius-bench-battery}) addresses this gap by
running every controller on both the observed and true state machines
simultaneously, reporting True-State Violation Rate (TSVR) alongside
Observed-State Violation Rate (OSVR).

\section{The Novelty Gap}

Table~\ref{tab:novelty-gap} summarises the positioning of DC3S relative
to the prior art.

\begin{table}[ht]
\centering
\caption{Positioning of DC3S relative to prior work.  A filled circle
(\ding{108}) indicates the capability is present; an open circle
(\ding{109}) indicates it is absent.}
\label{tab:novelty-gap}
\begin{tabular}{lccccc}
\toprule
\textbf{Method} &
\rotatebox{70}{\textbf{Obs.\ quality}} &
\rotatebox{70}{\textbf{Conformal UQ}} &
\rotatebox{70}{\textbf{True-state eval.}} &
\rotatebox{70}{\textbf{Runtime cert.}} &
\rotatebox{70}{\textbf{Adaptive width}} \\
\midrule
Deterministic LP   & \ding{109} & \ding{109} & \ding{109} & \ding{109} & \ding{109} \\
Robust MPC         & \ding{109} & \ding{109} & \ding{109} & \ding{109} & \ding{108} \\
Stochastic MPC     & \ding{109} & \ding{109} & \ding{109} & \ding{109} & \ding{108} \\
Classical shield   & \ding{109} & \ding{109} & \ding{109} & \ding{108} & \ding{109} \\
CBF shield         & \ding{109} & \ding{109} & \ding{109} & \ding{108} & \ding{109} \\
\textbf{DC3S (ours)} & \ding{108} & \ding{108} & \ding{108} & \ding{108} & \ding{108} \\
\bottomrule
\end{tabular}
\end{table}

\section{ORIUS as a Unifying Framework}
\label{sec:orius-unification}

The universality chapter (Chapter~\ref{ch:universality-completeness},
Theorem~\ref{thm:unification}) packages a precise unification claim:
CBF, Robust MPC, and Safe RL can each be represented as special cases of the
ORIUS degraded-observation kernel under fixed or model-derived reliability
policies.

\begin{description}
  \item[\textbf{CBF} (Control Barrier Functions).]
    CBF assumes $w_t \equiv 1$ (perfect observation).  The barrier constraint
    defines the tightened set; the CBF-QP is the repair step.  Under
    telemetry degradation, CBF does not inflate the margin, incurring the
    OASG phenomenon (T9 applies with large $c$).
  \item[\textbf{Robust MPC} (tube methods).]
    Robust MPC uses a fixed tube radius $r$, equivalent to $w_t \equiv
    1 - r/r_{\max}$ for all steps.  It applies a constant safety margin
    regardless of whether current telemetry is clean or degraded,
    sacrificing performance unnecessarily during clean periods.
  \item[\textbf{Safe RL} (posterior constraint).]
    Safe RL derives the action constraint from the posterior predictive
    distribution, which corresponds to $w_t$ from the posterior concentration
    and $q_t$ from the posterior quantile.  ORIUS provides the same guarantee
    without model assumptions via conformal prediction.
\end{description}

ORIUS strictly generalises each in the sense relevant to this dissertation:
dynamic $w_t$ from the OQE reduces unnecessary conservatism under clean
telemetry and automatically tightens under degraded telemetry.  The
reference reductions are implemented in \texttt{src/orius/universal/unification.py},
while the canonical production kernel remains
\texttt{src/orius/universal\_theory/}.

\section{Chapter Summary}

The literatures on battery dispatch, CPS telemetry faults, conformal
prediction, and runtime shielding each address a piece of the
observation-aware safety problem, but none provide a complete solution.
Dispatch optimisers ignore observation quality.  Fault-tolerant control
handles faults individually without feeding a unified score into the
safety layer.  Conformal methods provide coverage guarantees but do not
condition on telemetry reliability.  Runtime shields check constraints on
the observed state, missing the OASG entirely.  The ORIUS framework and
its DC3S realisation unify these strands into a single pipeline with
end-to-end safety guarantees from observation to actuation.
